\section{Постановка, две ключевые сцены и обозначения}
\emph{Прежде чем судить, на что система способна, проверьте, чего ей не дали узнать.}
Когда динамическая система может восстановить факт о собственном прошлом?
Сколько битов нужно ей передать и каким вмешательством можно проверить, что она
воспользовалась ими? Зафиксируем, какие величины доступны получателю к моменту
запроса. Две конструкции Rule~110 показывают, почему это существенно.

\paragraph{Два датчика покрывают окно и упускают различие.}
Поставим датчики в позициях $S=(-1,1)$ и прочтём их на тактах 0 и 1.
Нужно предсказать значение клетки 0 на такте 2. Всё окно её зависимости
$[-2,2]$ покрыто конусами датчиков. В первом начальном состоянии единица стоит только в $-1$,
во втором --- в $-1$ и $0$.
В обоих случаях датчики дважды показывают одну и ту же пару $(1,0)$.
На такте 2 целевая клетка равна соответственно 0 и 1. Какой бы ответ ни дал
декодер по этому журналу, в одном из двух случаев он ошибётся
(§\ref{sec:observer-coverage}).

\paragraph{Один бит возвращает утраченное различие.}
Теперь цель --- прежнее значение $Y$ клетки 5. В каждой из двух пар историй
переворот только $Y$ меняет показание $Z$ клетки 3 на такте 2. Наблюдатель
сохраняет этот бит. На такте 4 основные конфигурации внутри каждой пары
уже совпадают на всей прямой. Запишем $Z$ в клетку $-1$. Её левый сосед
равен 1, правый хранит контекстный бит $K$. Следующее обновление даёт
$Z\oplus K=Y$. Если при том же состоянии решётки вернуть $1-Z$, ответ
обратится во всех четырёх историях (§\ref{sec:observer-return}).

В этой конструкции по каналу идёт сырое показание датчика. В двухшаговом
опыте OLMo отправитель сначала по прошлому выбору $h$ и обратной связи $e$
вычисляет исправленное правило $p=h\oplus e$. В новый эпизод попадает только
выбранный до предъявления строки $\mathbf{x}$ токен $m\in\{\mathrm{A},\mathrm{B}\}$.
Получатель заново обрабатывает его в свежем KV-кэше. Так измеряется перенос
уже выбранного правила между эпизодами. Внутренняя коррекция стратегии при
непрерывном сохранении состояния требует отдельной проверки (§\ref{sec:olmo}).

\begin{center}\small
\begin{tabularx}{\linewidth}{@{}lX@{}}\toprule
Обозначение & Смысл на протяжении монографии \\\midrule
$\omega,H,W,Z$ & Допустимая история, сведения получателя, доступ наблюдателя и доставленное сообщение.\\
$Y$ & Целевой бит прошлого. В более общей постановке T1--T3 его алфавит может быть конечным.\\
$t,i,j,s$ & Дискретное время и целочисленные координаты решётки. $e_j$ меняет только начальную клетку $j$.\\
$B_t,\widetilde B_t,S_t$ & Невозмущённая и возмущённая конфигурации и множество координат, в которых они различаются.\\
$G_t,D_t,I_t$ & Значение клетки в начале координат, существенные начальные координаты и теоретическое окно зависимости.\\
$h,e,p,m$ & Прошлый выбор OLMo, бит обратной связи, исправленное правило и токен сообщения.\\
$\mathbf{x}=LwR$ & Новый восьмибитный вход. $w$ --- шестибитная середина, $L,R$ --- биты на концах.\\
$B,E(a,t)$ & Пакет мостовых постулатов и предикат опыта экземпляра $a$ в момент $t$.\\
\bottomrule\end{tabularx}
\end{center}

Формула указывает, что следует из заданных предпосылок (\texttt{COMPUTED}). Для вывода о данной машине нужен исполнимый свидетель или отчёт измерения (\texttt{MEASURED}). Статус \texttt{UNBOUND} означает отсутствие требуемой привязки, а \texttt{NOT\_RUN} --- отсутствие запуска в обозначенном срезе. Слово «точно» далее относится либо ко всем элементам заданной математической области, либо к явно указанной конечной панели. В §\ref{sec:observer-theory} сначала установим, при каких условиях точный декодер вообще существует.
